%% toc: P15 166 166.2
\chapter{Joint and Conditional Entropy, Dependence and Information Loss}
\label{sec:p15-w02b-p15-041}

Two chance variables observed together have a joint law, and each also
has its own law. Entropy measures the spread of each of these laws.
This chapter compares the three numbers and the conditional entropies
built from them. It then uses the comparison to measure how strongly the
variables depend on one another and how much is lost when an observation
is coarsened or used through a noisy channel.

\emph{Inputs.} A finite joint probability table, or a joint density
with respect to Lebesgue measure together with stated integrability
assumptions. \emph{Conclusions.} Exact identities and inequalities
among joint, marginal and conditional entropies; equality conditions
that identify independence or functional determination; a nonnegative
dependence measure; and, for densities, the rate of a pair and the way
each joint and conditional entropy of the pair moves when one variable
is recoded. The presentation follows Shannon's
corrected reprint of the 1948 paper: its Section~6 for finite laws
\cite[pp.~11--12]{Shannon1948MathematicalTheory}, its Section~12 for
equivocation and the rate of transmission
\cite[pp.~20--21]{Shannon1948MathematicalTheory}, and its Section~20
for pairs with a joint density
\cite[pp.~35--38]{Shannon1948MathematicalTheory}; the properties of a
single density and of a change of coordinates are proved in
\cref{sec:p15-w02b-p15-042-continuous-entropy}. The entropy of a
single finite distribution, its units and extreme values, and its
axiomatic characterization (Shannon's Theorem~2) are given in
\cref{sec:p15-w02b-p15-040-discrete}; see
\cref{def:p15-w02b-p15-040-entropy} and
\cref{thm:p15-w02b-p15-040-characterization}. Relative entropy for finite masses is used in
the form of \cref{thm:p15-rmdl-gibbs}.

\section{Finite joint laws and their entropies}
\label{sec:p15-w02b-p15-041-finite}

\begin{hbdefinition}{def:p15-w02b-p15-041-joint}{Joint, marginal and conditional entropies}
Let \(x\) take values in \(\{1,\ldots,m\}\) and \(y\) in
\(\{1,\ldots,n\}\), and let \(p(i,j)\geq0\), with
\(\sum_{i,j}p(i,j)=1\), be the probability that \(x=i\) and \(y=j\).
Write
\[
 p(i,\cdot)\defeq\sum_{j}p(i,j),\qquad p(\cdot,j)\defeq\sum_{i}p(i,j)
\]
for the two marginal masses. Fix a logarithm base \(b>1\); base \(2\)
gives bits. Every sum below runs only over cells with \(p(i,j)>0\),
equivalently \(0\log0=0\). On such a cell \(p(i,\cdot)\geq p(i,j)>0\)
and \(p(\cdot,j)>0\), so every logarithm written is finite.
Following Shannon \cite[p.~12]{Shannon1948MathematicalTheory}, define
\begin{align*}
 H(x,y)&\defeq-\sum_{i,j}p(i,j)\log p(i,j),\\
 H(x)&\defeq-\sum_{i,j}p(i,j)\log p(i,\cdot)
      =-\sum_{i}p(i,\cdot)\log p(i,\cdot),\\
 H(y)&\defeq-\sum_{i,j}p(i,j)\log p(\cdot,j)
      =-\sum_{j}p(\cdot,j)\log p(\cdot,j).
\end{align*}
For \(p(i,\cdot)>0\) the conditional probability that \(y=j\) given
\(x=i\) is \(p_i(j)\defeq p(i,j)/p(i,\cdot)\). The
\emph{conditional entropy of \(y\) given \(x\)} is
\[
 H_x(y)\defeq-\sum_{i,j}p(i,j)\log p_i(j).
\]
Exchanging the roles of the variables gives \(p_j(i)\defeq
p(i,j)/p(\cdot,j)\) and \(H_y(x)\defeq-\sum_{i,j}p(i,j)\log p_j(i)\).
For \(p(i,\cdot)>0\), write
\(H(y\mid x=i)\defeq-\sum_jp_i(j)\log p_i(j)\) for the entropy of the
\(i\)-th conditional law. The variables are \emph{independent} when
\(p(i,j)=p(i,\cdot)\,p(\cdot,j)\) for all \(i,j\).

In Shannon's notation \(x\) in \(H(x)\) labels a number; it is not
an argument of a function \cite[p.~11]{Shannon1948MathematicalTheory}.
A finite tuple of finite-valued variables with a common joint table is
again one finite-valued variable. Thus \(H(x,y,z)\), \(H_y(x,z)\) and
\(H_{yz}(x)\) are given by the same formulas with the tuple in place of
a single variable; \(H_{yz}(x)\) conditions on the pair \((y,z)\).
Changing the base multiplies every quantity in this definition by the
same positive constant, so each identity and inequality below holds in
every base.
\index{entropy!joint}\index{entropy!conditional}
\end{hbdefinition}

The source writes its sums over all \(i,j\). The positive-cell
convention is the reading under which a zero cell contributes nothing;
it also makes every displayed logarithm finite.

\begin{hbtheorem}{thm:p15-w02b-p15-041-average}{proved=annex:p15-w02b-p15-041-average}{Conditional entropy as an average}
Under \cref{def:p15-w02b-p15-041-joint},
\[
 H_x(y)=\sum_{i:\,p(i,\cdot)>0}p(i,\cdot)\,H(y\mid x=i).
\]
Hence \(H_x(y)\geq0\). Moreover \(H_x(y)=0\) exactly when, for each
\(i\) with \(p(i,\cdot)>0\), there is only one \(j\) with
\(p(i,j)>0\); that is, \(y\) is a function of \(x\) on the cells of
positive probability.
\end{hbtheorem}

The first identity is Shannon's description of \(H_x(y)\) as the
entropy of \(y\) for each value of \(x\), averaged with the weights of
those values \cite[p.~12]{Shannon1948MathematicalTheory}. The zero case
is a consequence of the averaged form.

\begin{hbtheorem}{thm:p15-w02b-p15-041-subadditivity}{proved=annex:p15-w02b-p15-041-subadditivity}{Subadditivity of joint entropy}
Under \cref{def:p15-w02b-p15-041-joint},
\[
 H(x,y)\leq H(x)+H(y),
\]
with equality if and only if \(x\) and \(y\) are independent.
\end{hbtheorem}

Shannon states the inequality with equality only under independence
\cite[p.~12]{Shannon1948MathematicalTheory}. The proof in the annex also
gives the converse: an independent pair attains equality.

\begin{hbtheorem}{thm:p15-w02b-p15-041-chain}{proved=annex:p15-w02b-p15-041-chain}{The chain identity and the conditioning bound}
Under \cref{def:p15-w02b-p15-041-joint},
\[
 H(x,y)=H(x)+H_x(y)=H(y)+H_y(x).
\]
Consequently
\[
 0\leq H_x(y)\leq H(y),
\]
and \(H_x(y)=H(y)\) if and only if \(x\) and \(y\) are independent.
\end{hbtheorem}

Shannon derives the first identity by substituting the conditional
probability and combines it with subadditivity to obtain
\(H(y)\geq H_x(y)\) \cite[p.~12]{Shannon1948MathematicalTheory}.
Knowing \(x\) never increases the average uncertainty about \(y\).
It lowers that uncertainty unless the two variables are independent.
This is an average statement. A single observed value \(x=i\) can have
\(H(y\mid x=i)>H(y)\); \cref{def:p15-w02b-p15-041-skew} gives a case.

\begin{hbtheorem}{thm:p15-w02b-p15-041-averaging}{proved=annex:p15-w02b-p15-041-averaging}{Averaging increases entropy}
Let \((p_1,\ldots,p_n)\) be a probability vector and let
\(a_{ij}\geq0\), \(1\leq i,j\leq n\), satisfy
\(\sum_ia_{ij}=\sum_ja_{ij}=1\). Put
\[
 p'_i=\sum_ja_{ij}p_j .
\]
Then \((p'_1,\ldots,p'_n)\) is a probability vector and
\(H(p'_1,\ldots,p'_n)\geq H(p_1,\ldots,p_n)\). Equality holds if and
only if \((p'_1,\ldots,p'_n)\) is a rearrangement of
\((p_1,\ldots,p_n)\).

In particular, if \(p_1<p_2\) and \(0<\delta<p_2-p_1\), replacing
\(p_1,p_2\) by \(p_1+\delta,p_2-\delta\) and leaving the other
coordinates unchanged strictly increases \(H\).
\end{hbtheorem}

Shannon states that such an averaging increases \(H\) except where the
transformation amounts to a permutation of the \(p_j\)
\cite[p.~12]{Shannon1948MathematicalTheory}. The theorem fixes the
reading of that exception. It concerns the averaged vector, not the
matrix. If all \(p_j\) are equal, every admissible matrix leaves the
vector unchanged, so the entropy is unchanged even when the matrix is
not a permutation matrix. The second paragraph is the source's
two-probability equalization, with the size of the transfer made
explicit.

\section{Dependence and information loss}
\label{sec:p15-w02b-p15-041-dependence}

\begin{hbdefinition}{def:p15-w02b-p15-041-rate}{Equivocation and the rate of transmission}
Under \cref{def:p15-w02b-p15-041-joint}, the \emph{equivocation} of
\(x\) given \(y\) is \(H_y(x)\), and
\[
 R\defeq H(x)-H_y(x).
\]
In Shannon's setting \(x\) is the transmitted symbol and \(y\) the
received symbol of a noisy channel, per symbol of text; \(H_y(x)\) is
the uncertainty about what was sent that remains on average after
reception, and \(R\) is the rate of actual transmission
\cite[p.~20]{Shannon1948MathematicalTheory}. The definition here applies
to any finite joint table. When the pair needs naming, write
\(R(x;y)\).
\index{equivocation}\index{rate of transmission}
\end{hbdefinition}

\begin{hbtheorem}{thm:p15-w02b-p15-041-rate}{proved=annex:p15-w02b-p15-041-rate}{Three forms of the rate and its range}
Under \cref{def:p15-w02b-p15-041-rate},
\[
 R=H(x)-H_y(x)=H(y)-H_x(y)=H(x)+H(y)-H(x,y).
\]
Hence \(R(x;y)=R(y;x)\) and
\[
 0\leq R\leq\min\{H(x),H(y)\}.
\]
Further, \(R=0\) if and only if \(x\) and \(y\) are independent, and
\(R=H(x)\) if and only if \(x\) is a function of \(y\) on the cells of
positive probability.
\end{hbtheorem}

The three forms are those Shannon records
\cite[p.~21]{Shannon1948MathematicalTheory}. The third shows that \(R\)
is the amount by which the joint entropy falls short of the sum of the
separate entropies. It is therefore a symmetric measure of statistical
dependence, zero exactly under independence by
\cref{thm:p15-w02b-p15-041-subadditivity}.

\begin{hbtheorem}{thm:p15-w02b-p15-041-correction}{proved=annex:p15-w02b-p15-041-correction}{Conditioning on further data}
Let \(x,y,z\) be finite-valued variables with a common joint table.
Then
\[
 H_y(x,z)=H_y(z)+H_{yz}(x)\geq H_y(x),
\]
\[
 H_{yz}(x)\geq H_y(x)-H_y(z)\geq H_y(x)-H(z),
\]
and \(H_{yz}(x)\leq H_y(x)\).
In particular, if \(H(z)<H_y(x)\), then \(H_{yz}(x)>0\).
\end{hbtheorem}

Shannon uses the first two lines in Section~12. There \(x\) is the
source output, \(y\) the received signal and \(z\) the data sent over a
correction channel. If the correction data carry less entropy than the
equivocation, uncertainty about what was sent remains even after both
signals are known \cite[p.~21]{Shannon1948MathematicalTheory}. The
theorem gives only this inequality step. The coding half of his
Theorem~10, which concerns long sequences and channel capacity, is a
separate result and is not proved here
\cite[pp.~20--21]{Shannon1948MathematicalTheory}.

\begin{hbtheorem}{thm:p15-w02b-p15-041-loss}{proved=annex:p15-w02b-p15-041-loss}{Coarsening and processing lose information}
Let \(x\) and \(y\) have a finite joint table.
\begin{enumerate}
\item For every function \(f\) from the values of \(x\) to a finite
set, \(H(f(x))\leq H(x)\). Equality holds if and only if \(f\) is
one-to-one on \(\{i:p(i,\cdot)>0\}\).
\item For every function \(g\) from the values of \(y\) to a finite
set, \(R(x;g(y))\leq R(x;y)\).
\end{enumerate}
\end{hbtheorem}

These two statements are consequences of the chain identity and of
\cref{thm:p15-w02b-p15-041-correction}; they are not stated in the
source. The first says that merging values of a variable cannot raise
its entropy and lowers it whenever two values of positive probability
are merged. The second says that no further computation on the
received signal alone can raise the rate of transmission it supports.

\begin{hbexample}{def:p15-w02b-p15-041-channel}{Shannon's binary channel and a comparison of dependence}
Use base \(2\). Let \(x\) take values \(0\) and \(1\) with probability
\(\tfrac12\) each, and let the received \(y\) differ from \(x\) with
probability \(0.01\), whatever was sent. The joint table is
\[
 p(0,0)=p(1,1)=0.495,\qquad p(0,1)=p(1,0)=0.005 .
\]
Both marginals are \((\tfrac12,\tfrac12)\), so \(H(x)=H(y)=1\) bit.
Given either received value, the sent value is correct with
probability \(0.99\). By \cref{thm:p15-w02b-p15-041-average},
\[
 H_y(x)=-(0.99\log_20.99+0.01\log_20.01)\approx0.0808
\]
bits per symbol, and \(R=1-0.0808\approx0.919\). Shannon reports
\(0.081\) and \(919\) bits per second at \(1000\) symbols per second
\cite[p.~20]{Shannon1948MathematicalTheory}. As a check on the third
form, \(H(x,y)=1+0.0808\), so \(H(x)+H(y)-H(x,y)=0.919\).

Replace the error probability by \(\tfrac12\). Then
\(p(i,j)=\tfrac14\) for all four cells, the variables are independent,
\(H_y(x)=1\) and \(R=0\). Half of the received symbols still agree with
the sent ones, but by chance alone. Counting correct symbols would
credit this channel with transmission it does not perform; the rate
does not \cite[p.~20]{Shannon1948MathematicalTheory}.

Finally, process the received symbol of the first channel by the constant
map \(g(y)=0\). Then \(g(y)\) is independent of \(x\), and
\(R(x;g(y))=0<R(x;y)\), as the processing bound allows.
\end{hbexample}

\begin{hbexample}{def:p15-w02b-p15-041-skew}{One observed value can raise uncertainty}
Use base \(2\). Let \(p(1,1)=0.9\), \(p(2,1)=p(2,2)=0.05\), and
\(p(1,2)=0\). Then \(p(\cdot,1)=0.95\) and
\(H(y)=-(0.95\log_20.95+0.05\log_20.05)\approx0.286\). Given \(x=2\),
the law of \(y\) is \((\tfrac12,\tfrac12)\), so
\(H(y\mid x=2)=1>H(y)\). Given \(x=1\), \(y=1\) for certain. The
average is \(H_x(y)=0.1\cdot1=0.1<0.286\), in accordance with
\cref{thm:p15-w02b-p15-041-chain}.
\end{hbexample}

\section{Densities: dependence and recoding}
\label{sec:p15-w02b-p15-041-continuous}

For a continuous distribution Shannon defines entropy by the integral
analogue of the finite sum, together with joint and conditional forms
\cite[p.~35]{Shannon1948MathematicalTheory}. Those definitions, the
conditions under which the integrals exist, and the continuous
counterparts of the finite results of
\cref{sec:p15-w02b-p15-041-finite} are given and proved in
\cref{sec:p15-w02b-p15-042-continuous-entropy}: the volume bound
(\cref{thm:p15-w02b-p15-042-volume}), subadditivity
(\cref{thm:p15-w02b-p15-042-joint}), averaging
(\cref{thm:p15-w02b-p15-042-averaging}), the chain identity and the
conditioning bound (\cref{thm:p15-w02b-p15-042-conditional}), the
Gaussian and exponential maxima
(\cref{thm:p15-w02b-p15-042-gaussian,thm:p15-w02b-p15-042-exponential})
and the change of entropy under a change of coordinates
(\cref{thm:p15-w02b-p15-042-coordinate-change-thm,thm:p15-w02b-p15-042-linear-change}).
This section adds what concerns dependence: the rate for a pair with a
joint density, which finite properties it keeps and loses, and how each
joint and conditional entropy of a pair moves when one variable is
recoded.

\begin{hbdefinition}{def:p15-w02b-p15-041-density}{The rate for a pair with a joint density}
Let \(x\) and \(u\) take values in \(\mathbb R^n\) and \(\mathbb R^k\)
and have a joint density \(p(x,u)\) with respect to Lebesgue measure on
\(\mathbb R^{n+k}\), with marginal densities \(p(x)=\int p(x,u)\,du\)
and \(p(u)=\int p(x,u)\,dx\). The entropies \(H(x)\), \(H(u)\),
\(H(x,u)\) and the conditional entropies \(H_x(u)\), \(H_u(x)\) are those
of \cref{def:p15-w02b-p15-042-reference-entropy} in Shannon's case, with
Lebesgue measure as reference measure and natural logarithms. Say that the pair has
\emph{finite entropies} when \(H(x)\), \(H(u)\) and \(H(x,u)\) are
finite real numbers. Then \(H_x(u)\) and \(H_u(x)\) are finite by
\cref{thm:p15-w02b-p15-042-conditional}, and the \emph{rate} of the
pair is
\[
 R(x;u)\defeq H(x)-H_u(x).
\]
The variables are \emph{independent} when \(p(x,u)=p(x)p(u)\) for almost
every \((x,u)\). As in the source, \(p(x)\) and \(p(u)\) are two
different functions distinguished by the names of their arguments
\cite[p.~35]{Shannon1948MathematicalTheory}.
\index{rate of transmission!for densities}
\end{hbdefinition}

\begin{hbtheorem}{thm:p15-w02b-p15-041-density-rate}{proved=annex:p15-w02b-p15-041-density-rate}{The rate of a pair with a joint density}
Let the pair \((x,u)\) of \cref{def:p15-w02b-p15-041-density} have
finite entropies. Then
\[
 R(x;u)=H(x)-H_u(x)=H(u)-H_x(u)=H(x)+H(u)-H(x,u),
\]
so \(R(x;u)=R(u;x)\). Moreover \(R(x;u)\geq0\), with equality if and only
if \(x\) and \(u\) are independent.
\end{hbtheorem}

This is the continuous form of \cref{thm:p15-w02b-p15-041-rate}. Two
parts of the finite statement do not carry over. The conditional
entropy \(H_u(x)\) can be negative
(\cref{thm:p15-w02b-p15-042-conditional}), and the bound
\(R\leq\min\{H(x),H(u)\}\) fails: in
\cref{def:p15-w02b-p15-041-recoding-example} the rate is positive while
both marginal entropies are negative. Shannon's remark that rates are
nonnegative although continuous entropies can be negative
\cite[p.~38]{Shannon1948MathematicalTheory} is the first part of the
inequality above.

\begin{hbtheorem}{thm:p15-w02b-p15-041-coordinates}{proved=annex:p15-w02b-p15-041-coordinates}{Recoding one variable of a pair}
Let \((T,j)\) be a coordinate change from \(U\) to \(V\) in
\(\mathbb R^n\) in the sense of
\cref{def:p15-w02b-p15-042-coordinate-change}, and let \(x\) have a
density \(p(x)\) with respect to Lebesgue measure, with
\(P(x\in U)=1\), \(H(x)\) finite and
\(\int p(x)\,|\log j(x)|\,dx<\infty\). Put \(y=T(x)\), and write
\(|J(x/y)|=1/j(x)\) for the absolute value of Shannon's Jacobian of the
old coordinates with respect to the new.
\begin{enumerate}
\item The entropy \(H(y)\) is finite and
\[
 H(y)=H(x)+\int p(x)\log j(x)\,dx=H(x)-\int p(x)\log|J(x/y)|\,dx .
\]
\item Let, in addition, \(u\) be a random vector in \(\mathbb R^k\) such
that the pair \((x,u)\) satisfies \cref{def:p15-w02b-p15-041-density}
with finite entropies. Then the pair \((y,u)\) also has a joint density
and finite entropies, and, with \(c=\int p(x)\log j(x)\,dx\),
\[
 H(y,u)=H(x,u)+c,\qquad H_u(y)=H_u(x)+c,\qquad H_y(u)=H_x(u),
\]
\[
 R(y;u)=R(x;u).
\]
\end{enumerate}
\end{hbtheorem}

Item~1 is Shannon's formula: the new entropy is the old one less the
expected logarithm of the Jacobian
\cite[pp.~37--38]{Shannon1948MathematicalTheory}; it restates
\cref{thm:p15-w02b-p15-042-coordinate-change-thm}. Item~2 sorts the
quantities of a pair by their behaviour. Every entropy in which the
recoded variable is the one whose spread is measured shifts by the same
constant \(c\). An entropy of \(u\), conditional on \(y\) or not, does not
move, because conditioning on \(y\) and conditioning on \(x\) are
conditioning on the same information. The rate is a difference of two
terms that shift together, which is the precise content of Shannon's
remark that rates and capacities do not depend on the coordinate frame
\cite[p.~38]{Shannon1948MathematicalTheory}. When both variables are
recoded separately, \cref{thm:p15-w02b-p15-042-difference-invariance}
gives the same conclusion for the rate.

\begin{hbexample}{def:p15-w02b-p15-041-recoding-example}{A Gaussian pair recoded three ways}
Use natural logarithms. Let \(z_1,z_2\) be independent standard normal
variables, fix \(\sigma>0\), and put
\[
 x=\sigma z_1,\qquad u=\sigma\bigl(0.6\,z_1+0.8\,z_2\bigr).
\]
This is the pair of \cref{thm:p15-w02b-p15-042-boundaries}, item~4,
with \(\sigma_1=\sigma_2=\sigma\) and \(\rho=0.6\), so
\(\sqrt{1-\rho^2}=0.8\). Each of \(x\) and \(u\) is normal with mean
\(0\) and variance \(\sigma^2\), and by that result and
\cref{thm:p15-w02b-p15-042-gaussian}, item~5,
\[
 H(x)=H(u)=\log\bigl(\sqrt{2\pi e}\,\sigma\bigr),\qquad
 H(x,u)=\log\bigl(2\pi e\,\sigma^2\cdot0.8\bigr),
\]
\[
 R(x;u)=-\tfrac12\log0.64=\log1.25\approx0.2231 .
\]
By \cref{thm:p15-w02b-p15-041-density-rate},
\(H_x(u)=H(x,u)-H(x)=\log\bigl(\sqrt{2\pi e}\,\sigma\cdot0.8\bigr)\).

\emph{Negative entropies, positive rate.} Take \(\sigma=0.1\). Since
\(\tfrac12\log(2\pi e)\approx1.4189\) and \(\log0.1\approx-2.3026\),
\(H(x)=H(u)\approx-0.8837\), while \(R(x;u)\approx0.2231\). The rate
exceeds both marginal entropies, which no finite table allows
(\cref{thm:p15-w02b-p15-041-rate}).

\emph{A linear recoding of one variable.} Let \(y=2x\), a coordinate
change of \(\mathbb R\) with constant factor \(j=2\)
(\cref{def:p15-w02b-p15-042-coordinate-change}, item~1). By
\cref{thm:p15-w02b-p15-041-coordinates}, \(H(y)=H(x)+\log2\),
\(H(y,u)=H(x,u)+\log2\) and \(H_u(y)=H_u(x)+\log2\), while
\(H_y(u)=H_x(u)\) and \(R(y;u)=\log1.25\) as before.

\emph{A nonlinear recoding with no net shift.} Let \(y=e^{x}\), a
continuously differentiable map of \(\mathbb R\) onto \((0,\infty)\)
with positive derivative \(j(x)=e^{x}\)
(\cref{def:p15-w02b-p15-042-coordinate-change}, item~2). Here
\(\int p(x)\,|\log j(x)|\,dx\) is the mean of \(|x|\), which is finite
because \(|x|\leq1+x^2\) and \(x\) has finite variance; and
\(c=\int p(x)\,x\,dx=0\). Hence \(H(e^{x})=H(x)\) and every quantity of
item~2 is unchanged, although the map stretches some intervals and
shrinks others. Equal entropies do not show that a recoding preserves
volume; only the average of \(\log j\) under the law of \(x\) enters.

\emph{A recoding that mixes the variables.} The theorem concerns
recoding one variable while the other stays fixed. The map
\((z_1,z_2)\mapsto(x,u)\) is a linear coordinate change of the plane with
determinant \(0.8\,\sigma^2\), so by
\cref{thm:p15-w02b-p15-042-linear-change}
\(H(x,u)=H(z_1,z_2)+\log(0.8\,\sigma^2)\). But it mixes the two
coordinates: \(z_1\) and \(z_2\) are independent, so
\(R(z_1;z_2)=0\) by \cref{thm:p15-w02b-p15-041-density-rate}, while
\(R(x;u)=\log1.25>0\). Dependence between two variables is a property of
the pair of coordinate systems, one for each variable; a joint
recoding can create or remove it.
\end{hbexample}

\section{Scope, relationships and sources}
\label{sec:p15-w02b-p15-041-scope}

\emph{What the quantities are not.} Every quantity in this chapter is a
functional of a probability law, together with a reference measure in
the continuous case. Entropy and conditional entropy measure spread of
that law; they say nothing about whether a message is true
(\cref{def:p15-w02b-p15-040-not-truth} computes contrasts). The rate
\(R\) measures statistical dependence of the joint law. It is not a
measure of evidential support for a hypothesis, which requires an
explicit hypothesis space and likelihoods, and it is not a decision
value, which requires actions and utilities. A positive \(R\) also
does not identify a direction of influence: \(R(x;y)=R(y;x)\).

\emph{Assumptions that limit the results.} The finite results need a
complete joint table. The continuous results need a joint density with
respect to Lebesgue measure and finite entropies; without them the
identities can become \(\infty-\infty\)
(\cref{thm:p15-w02b-p15-042-defined-cases}). The recoding result needs a
coordinate change in the sense of
\cref{def:p15-w02b-p15-042-coordinate-change}; for a continuously
differentiable map of several variables that is not linear, that
hypothesis is the multivariable change-of-variables theorem, which the
handbook does not prove and which must be verified for the map at hand.
Shannon's Theorem~10, which interprets the equivocation as the capacity
a correction channel must have, concerns long sequences; only its
inequality step is proved here
(\cref{thm:p15-w02b-p15-041-correction}).

\emph{Relationships.} The nonnegativity of relative entropy used in the
finite proofs is \cref{thm:p15-rmdl-gibbs}. The entropy of a single
finite distribution, its characterization and computed contrasts with
truth, support and value are in \cref{sec:p15-w02b-p15-040-discrete};
a further comparison of two tables with equal marginals is
\cref{def:p15-w02b-p15-040-dependence}. The finite-alphabet entropy of
\cref{def:p15-finite-alphabet-source} is \(H(x)\) for a single symbol,
and the asymptotic equipartition result there concerns sequences rather
than pairs. The differential entropy used for measurement models in
\cref{sec:p15-measurement-input-evaluation} is the density entropy of
\cref{def:p15-w02b-p15-042-reference-entropy} in natural logarithms.
The minimal-change updates of \cref{thm:p15-rat41-minimal-update} use
the relative entropy whose nonnegativity drives
\cref{thm:p15-w02b-p15-041-subadditivity}.

\emph{Sources.} Shannon's corrected reprint: Section~6, properties 3
to 6 \cite[p.~12]{Shannon1948MathematicalTheory}; Section~12, the
equivocation, the rate, Theorem~10 and the three forms of the rate
\cite[pp.~20--21]{Shannon1948MathematicalTheory}; Section~20, the
joint and conditional definitions and the coordinate remarks
\cite[pp.~35--38]{Shannon1948MathematicalTheory}.
\index{information!rate of transmission}
